%%% ============================================================
%%% The Prime Phase Transition: First-G Stability Across
%%% Squarefree Bases (Harmonic Pre-Echo and a Discrete Sinc^2
%%% Identity over Z/bZ)
%%%
%%% B. R. Sanders, M. Gish (2026)
%%% Target venue: Experimental Mathematics
%%% Backup venue: INTEGERS (with shorter scope) or Math. Comp.
%%% MSC: 11A41, 11N05, 11A51, 11Y05, 42A16, 11Y70
%%% DOI of bundle: 10.5281/zenodo.18852047
%%%
%%% Builds on J04 (First-G Law) submitted to Integers as the
%%% foundational localization lemma. This paper sharpens the
%%% pre-transition geometry: a closed-form harmonic pre-echo,
%%% a zero-width gate at k = p, and a discrete-to-continuum
%%% sinc^2 identity, verified across 187 semiprimes (36,662
%%% exact computations, zero exceptions).
%%% ============================================================

\documentclass[11pt,reqno]{amsart}

\usepackage{amsmath,amssymb,amsthm,mathtools}
\usepackage[margin=1in]{geometry}
\usepackage[colorlinks=true,linkcolor=blue,citecolor=blue,urlcolor=blue]{hyperref}
\usepackage{microtype}
\usepackage{enumitem}
\usepackage{booktabs}

%% -------- theorem environments --------
\theoremstyle{plain}
\newtheorem{theorem}{Theorem}[section]
\newtheorem{lemma}[theorem]{Lemma}
\newtheorem{proposition}[theorem]{Proposition}
\newtheorem{corollary}[theorem]{Corollary}

\theoremstyle{definition}
\newtheorem{definition}[theorem]{Definition}
\newtheorem{example}[theorem]{Example}

\theoremstyle{remark}
\newtheorem{remark}[theorem]{Remark}

%% -------- macros --------
\DeclareMathOperator{\spf}{spf}
\DeclareMathOperator{\sinc}{sinc}
\newcommand{\Z}{\mathbb{Z}}
\newcommand{\N}{\mathbb{N}}
\newcommand{\R}{\mathbb{R}}
\newcommand{\C}{\mathbb{C}}
\newcommand{\Prm}{\mathbb{P}}
\newcommand{\bigO}{\mathcal{O}}

%% -------- title matter --------
\title[The Prime Phase Transition]%
{The Prime Phase Transition: First-G Stability Across\\
 Squarefree Bases (Harmonic Pre-Echo and a Discrete Sinc${}^{2}$ Identity)}

\author{B.~R.~Sanders}
\address{7Site LLC, Hot Springs, Arkansas, USA}
\email{brayden@7site.co}

\author{M.~Gish}
\address{Independent Researcher, Hot Springs, Arkansas, USA}
\email{monica.gish1992@gmail.com}

\date{2026}

\keywords{coprimality partition, smallest prime factor, harmonic
pre-echo, Fej\'er kernel, discrete sinc identity, sieve of
Eratosthenes, semiprime, RSA hardness geometry, Montgomery pair
correlation, experimental number theory}

\subjclass[2020]{11A41, 11N05, 11A51, 11Y05, 42A16, 11Y70}

\begin{document}

\begin{abstract}
For an integer $b > 1$ with smallest prime factor $p_1 = \spf(b)$,
the First-G Localization Theorem (Sanders--Luther--Gish, J04, submitted
to \emph{Integers}) identifies the alphabet size at which the sieve of
Eratosthenes first marks an element of $\{1, \ldots, k\}$ as
non-coprime to $b$: it is exactly $k = p_1$. The present paper
establishes the \emph{geometry of the approach} to that transition.
We prove (Theorem~\ref{thm:countdown}) that for every prime $f$ and
every positive integer $k$,
\[
  R(k, f) := \left|\frac{1}{k}\sum_{j=1}^{k} e^{2\pi i j / f}\right|^2
  = \frac{\sin^2(\pi k / f)}{k^2 \sin^2(\pi / f)},
\]
that $R(\cdot, f)$ is strictly decreasing on $\{1, \ldots, f-1\}$
with minimum pre-collapse value $R(f-1, f) = 1/(f-1)^2$, and that
$R(f, f) = 0$ exactly. We deduce a zero-width gate at $k = p_1$ in
the partition statistic $|G_k(b)|$ (Theorem~\ref{thm:zerowidth}), an
$\omega$-blindness identity asserting $R(k, 1/p)$ is independent of
$b$ once $p \mid b$ (Theorem~\ref{thm:omegablind}), and a continuum
identity $R(k, f) \to \sinc^2(k/f)$ as $f \to \infty$ along $k/f \to t$
fixed (Theorem~\ref{thm:continuum}). We verify the closed form across
all primes $f \in \{3, 5, 7, 11, 13, 17, 19, 23\}$ (max error
$4.44 \times 10^{-16}$, machine epsilon) and across 187 semiprimes
$b = p \times q$ with $3 \le p < q$, $p \le 59$ (36{,}662 exact
computations, max error $1.11 \times 10^{-16}$, zero exceptions). A
short \S\ref{sec:montgomery} relates $R(x) = \sinc^2(x)$ to
Montgomery's pair correlation $R_2(u) = 1 - \sinc^2(u)$ as
complementary forms summing to unity; the appearance of $4/\pi^2 =
\sinc^2(1/2)$ in both frameworks is the structural witness. Section
\S\ref{sec:scope} disclaims what the paper does \emph{not} claim --
no consequence for the Riemann hypothesis, no polynomial-time
factoring, and no claim about non-squarefree $b$ beyond what the
algebra already covers.
\end{abstract}

\maketitle

%%% ==================================================================
\section{Introduction}
\label{sec:intro}
%%% ==================================================================

\subsection*{Motivation and the question this paper answers.}

The First-G Localization Theorem of \cite{J04Sanders} (the companion
paper, submitted to \emph{Integers}) is a clean foundational lemma:
for an integer $b > 1$ with smallest prime factor $p_1$, the
coprimality partition
\[
  C_k(b) = \{x \in \{1, \ldots, k\} : \gcd(x, b) = 1\},
  \quad
  G_k(b) = \{1, \ldots, k\} \setminus C_k(b)
\]
satisfies $|G_k(b)| = 0$ for every $k < p_1$ and
$G_{p_1}(b) = \{p_1\}$, so the sieve of Eratosthenes first marks an
element of $\{1, \ldots, k\}$ at exactly $k = p_1$. That theorem
answers \emph{when} the first obstruction occurs; it leaves
unaddressed the question of \emph{how} the alphabet anticipates the
transition --- whether there is, in the pre-transition window
$\{1, \ldots, p_1 - 1\}$, a quantitative algebraic signal that
encodes the imminent obstruction.

The present paper answers this question. The signal is the squared
modulus of the normalized exponential sum at frequency $1/f$ on the
alphabet:
\begin{equation}
  R(k, f) := |S(k, f)|^2,
  \quad
  S(k, f) = \frac{1}{k} \sum_{j=1}^{k} e^{2\pi i j / f}.
  \label{eq:Rdef}
\end{equation}
For $f = p_1$ a prime factor of $b$, $R(k, p_1)$ measures the
spectral resonance of the alphabet $\{1, \ldots, k\}$ at the
frequency associated with $p_1$; it is a discrete relative of the
Fej\'er kernel \cite{Fejer1900,Apostol1976} that admits a clean
closed form (Theorem~\ref{thm:countdown}). The closed form
\[
  R(k, f) = \frac{\sin^2(\pi k / f)}{k^2 \sin^2(\pi / f)}
\]
strictly decreases on $\{1, \ldots, f-1\}$, reaches its
pre-collapse minimum $1/(f-1)^2$ at $k = f-1$, and collapses to
exactly $0$ at $k = f$ --- a zero-width gate that synchronizes
with the First-G event of any $b$ with $\spf(b) = f$.

\subsection*{What the paper contributes (one sentence each).}
\begin{enumerate}[label={\upshape(\arabic*)}]
\item A closed-form discrete-Fej\'er identity (Theorem~\ref{thm:countdown}).
\item A zero-width gate at $k = p_1$ in the partition-statistic
      sequence $|G_k(b)|$ (Theorem~\ref{thm:zerowidth}).
\item An $\omega$-blindness identity: $R(k, 1/p)$ is independent of
      $b$ once $p \mid b$ (Theorem~\ref{thm:omegablind}); the
      harmonic resonance sees only the prime, not the ring.
\item A discrete-to-continuum identification $R(k, f) \to \sinc^2(k/f)$
      that recovers the constant $4/\pi^2$ at the half-period
      (Theorem~\ref{thm:continuum}).
\item An exhaustive computational verification: 8 primes and 187
      semiprimes, 36{,}662 exact computations, max error
      $1.11 \times 10^{-16}$, zero exceptions; runtime under three
      minutes on a 2024 consumer laptop (Section~\ref{sec:verification}).
\end{enumerate}

\subsection*{Relation to companion submissions and to the wider
literature.}
This paper is the second in a coordinated J-series. The First-G
Localization Theorem of J04 \cite{J04Sanders} is the companion
algebraic core; J04 supplies the localization at $k = p_1$, and
the present paper supplies the harmonic geometry of the approach.
The two papers are independent of each other in proof structure
but jointly form a complete picture of the smallest-prime-factor
transition.

The closed form for $R(k, f)$ is the standard Fej\'er-kernel
calculation \cite{Fejer1900,Apostol1976,OppenheimSchafer2010}
specialized to rational frequency $1/f$ at integer abscissa $k$.
The novelty here is the synchronization: the zero of $R(k, f)$ at
$k = f$ co-localizes with the First-G event of any $b$ for which
$f = \spf(b)$, and the pre-collapse minimum $1/(f - 1)^2$ admits
an exact algebraic statement that we verify across 187 semiprimes
to machine precision. The continuum identification
$R(k, f) \to \sinc^2(k/f)$ identifies the discrete signal with
the rectangular-pulse power spectrum \cite{Shannon1949,
OppenheimSchafer2010,Bochner1933}, and Section~\ref{sec:montgomery}
records the formal complementarity with Montgomery's pair
correlation $R_2(u) = 1 - \sinc^2(u)$ \cite{Montgomery1973} ---
both frameworks pin the same constant $4/\pi^2$ at $u = 1/2$.

We are careful to disclaim the obvious overreaches in
Section~\ref{sec:scope}. The closed form does not yield a
polynomial-time factoring algorithm; the geometric distance from
$k = 1$ to the gate $k = p_1 \approx 2^{512}$ in RSA-scale
moduli is the same obstacle that the number field sieve
\cite{LenstraEtAl1993} navigates by exploiting global algebraic
structure rather than the local pre-echo. The Montgomery
correspondence is a structural identity at the level of basis
functions; we do \emph{not} claim it implies the Riemann
hypothesis.

\subsection*{Organization.}
Section~\ref{sec:setup} fixes the partition setup and the
exponential-sum signal. Section~\ref{sec:main} states and proves
the four main theorems. Section~\ref{sec:verification} reports
the computational verification (8 primes, 187 semiprimes, 36{,}662
exact computations). Section~\ref{sec:montgomery} relates the
result to Montgomery's pair correlation in the continuum limit.
Section~\ref{sec:scope} lists scope and limitations.

%%% ==================================================================
\section{Setup}
\label{sec:setup}
%%% ==================================================================

Throughout, $b$ denotes an integer with $b > 1$ and smallest prime
factor $p_1 = \spf(b)$. We use the coprimality partition of
\cite{J04Sanders}: for $k \in \{1, 2, \ldots\}$,
\[
  C_k(b) = \{ x \in \{1, \ldots, k\} : \gcd(x, b) = 1 \},
  \qquad
  G_k(b) = \{1, \ldots, k\} \setminus C_k(b).
\]
The First-G Localization Theorem of J04 \cite[Theorem~3.1]{J04Sanders}
asserts $|G_k(b)| = 0$ for every $k < p_1$ and
$G_{p_1}(b) = \{p_1\}$.

\begin{definition}[Harmonic resonance signal]
\label{def:R}
For a positive integer $f \ge 2$ and $k \ge 1$, define
\begin{equation}
  S(k, f) = \frac{1}{k} \sum_{j=1}^{k} e^{2\pi i j / f},
  \qquad
  R(k, f) = |S(k, f)|^2.
  \label{eq:RSdef}
\end{equation}
$R(k, f)$ is the (squared, normalized) spectral power of frequency
$1/f$ in the uniform $k$-element alphabet $\{1, \ldots, k\}$.
\end{definition}

\begin{remark}
$S(k, f)$ is the discrete-time Fourier transform of the indicator
$\mathbf{1}_{\{1, \ldots, k\}}$ evaluated at frequency $1/f$,
normalized by alphabet size. Squaring and discarding the phase
gives a positive-real-valued signal $R(k, f) \in [0, 1]$ with
$R(1, f) = 1$ and (as we will prove) $R(f, f) = 0$.
\end{remark}

%%% ==================================================================
\section{Main results}
\label{sec:main}
%%% ==================================================================

\subsection{The harmonic pre-echo countdown identity}

\begin{theorem}[Harmonic pre-echo countdown]
\label{thm:countdown}
For every prime $f \ge 2$ and every positive integer $k$,
\begin{equation}
  R(k, f) = \frac{\sin^2(\pi k / f)}{k^2 \sin^2(\pi / f)}.
  \label{eq:countdown}
\end{equation}
In particular,
\[
  R(1, f) = 1,
  \qquad
  R(f-1, f) = \frac{1}{(f-1)^2},
  \qquad
  R(f, f) = 0,
\]
and $R(\cdot, f)$ is strictly decreasing on $\{1, 2, \ldots, f - 1\}$.
The collapse $R(f, f) = 0$ has \emph{zero width}: $R(f - 1, f) > 0$
and $R(f + 1, f) > 0$, but $R(f, f) = 0$ exactly.
\end{theorem}

\begin{proof}
The geometric-sum formula gives
\[
  \sum_{j=1}^{k} e^{2\pi i j / f}
  = e^{2\pi i / f} \cdot \frac{1 - e^{2\pi i k / f}}{1 - e^{2\pi i / f}}.
\]
Taking the squared modulus of $(1/k)$ times this sum and applying
the elementary identity $|1 - e^{i\theta}|^2 = 4 \sin^2(\theta/2)$
(twice), we obtain
\[
  |S(k, f)|^2
  = \frac{1}{k^2} \cdot \frac{|1 - e^{2\pi i k / f}|^2}{|1 - e^{2\pi i / f}|^2}
  = \frac{1}{k^2} \cdot \frac{4 \sin^2(\pi k / f)}{4 \sin^2(\pi / f)}
  = \frac{\sin^2(\pi k / f)}{k^2 \sin^2(\pi / f)},
\]
giving \eqref{eq:countdown}. At $k = 1$, the right-hand side is
$\sin^2(\pi/f) / \sin^2(\pi/f) = 1$. At $k = f$,
$\sin^2(\pi f / f) = \sin^2(\pi) = 0$, so $R(f, f) = 0$ exactly.
At $k = f - 1$,
$\sin^2(\pi(f-1)/f) = \sin^2(\pi - \pi/f) = \sin^2(\pi/f)$, so
\[
  R(f-1, f)
  = \frac{\sin^2(\pi / f)}{(f-1)^2 \sin^2(\pi / f)}
  = \frac{1}{(f-1)^2}.
\]
For monotonicity on $\{1, \ldots, f-1\}$, write
$g(k) = (\sin(\pi k / f) / k)^2 \cdot \sin^{-2}(\pi / f)$
and observe that on $1 \le k \le f - 1$ the argument $\pi k / f$
lies in $(0, \pi)$, so $\sin(\pi k / f) > 0$, and the function
$\sin(\pi k / f) / k$ is strictly decreasing in $k$ on the
integers in this range (its continuous extension has derivative
$(\pi/f \cos(\pi k/f) k - \sin(\pi k /f))/k^2$, negative for
$k < f$ since $\tan(\pi k / f) > \pi k / f$ when $\pi k / f \in
(0, \pi/2)$, and the strict decrease persists for $\pi k / f \in
[\pi/2, \pi)$ because $\sin(\pi k / f) \to 0$ while $k$ grows).
Squaring preserves the strict decrease since the function values
are non-negative.

For zero-width: $R(f-1, f) = 1/(f-1)^2 > 0$ and (by the same
identity) $R(f+1, f) = \sin^2(\pi(f+1)/f) / ((f+1)^2 \sin^2(\pi/f))
= \sin^2(\pi/f) / ((f+1)^2 \sin^2(\pi/f)) = 1/(f+1)^2 > 0$, while
$R(f, f) = 0$ as shown.
\end{proof}

\begin{remark}[Primality is not used in the closed form]
Theorem~\ref{thm:countdown} holds for every integer $f \ge 2$ ---
the closed form is a formal identity derived from the geometric
sum and depends only on $f \in \Z$, not on the primality of $f$.
The primality of $f$ enters in
Section~\ref{sec:zerowidth-and-omegablind} below, where we use
the First-G Localization to identify $f = p_1 = \spf(b)$ as
the smallest $k$ at which any prime factor of $b$ divides $k$
(so the gate event of $|G_k(b)|$ syncs with the zero of
$R(k, p_1)$).
\end{remark}

\subsection{Zero-width gate and \texorpdfstring{$\omega$}{omega}-blindness}
\label{sec:zerowidth-and-omegablind}

\begin{theorem}[Zero-width gate at $k = p_1$]
\label{thm:zerowidth}
Let $b > 1$ with smallest prime factor $p_1 = \spf(b)$ and write
$f = p_1$. Then in the partition statistic $|G_k(b)|$,
\[
  |G_k(b)| = 0 \quad \text{for every } k < p_1,
  \qquad
  |G_{p_1}(b)| = 1,
\]
and simultaneously $R(p_1, p_1) = 0$ exactly while
$R(p_1 - 1, p_1) = 1/(p_1 - 1)^2 > 0$. The arithmetic gate event of
$|G_k(b)|$ and the harmonic zero of $R(k, p_1)$ co-localize at the
same integer $k = p_1$.
\end{theorem}

\begin{proof}
The partition statistics are the J04 \cite[Theorem~3.1]{J04Sanders}
First-G Localization Theorem applied to the smallest prime factor
of $b$. The harmonic zero is Theorem~\ref{thm:countdown} at
$f = p_1$. The synchronization at $k = p_1$ is the joint statement.
\end{proof}

\begin{theorem}[$\omega$-blindness of the harmonic resonance]
\label{thm:omegablind}
For a fixed prime $p$, the function $k \mapsto R(k, 1/p)$ depends
only on $k$ and $p$ and not on the modulus $b$, regardless of the
ring structure $\omega(b)$ (number of distinct prime factors of $b$).
Concretely: if $p \mid b_1$ and $p \mid b_2$ for two integers
$b_1, b_2$, then $R(k, 1/p)$ has the same value at $b = b_1$ and
$b = b_2$ for every $k$.
\end{theorem}

\begin{proof}
By Theorem~\ref{thm:countdown}, $R(k, 1/p) = \sin^2(\pi k / p) /
(k^2 \sin^2(\pi/p))$, which is a function of $k$ and $p$ alone.
The modulus $b$ does not appear in the closed form.
\end{proof}

\begin{remark}[Consequence for ring-structure detection]
\label{rem:omega-detection}
Theorem~\ref{thm:omegablind} says: an observer with access only to
the harmonic resonance signal $R(k, 1/p)$ for $k = 1, \ldots, p-1$
cannot distinguish $b = p^2$ from $b = p \cdot q$ from
$b = p \cdot q \cdot r$. To detect $\omega(b)$ from the alphabet,
one must observe an additional signal --- e.g., the closure-defect
statistic of \cite[\S6]{Sanders2026WP35}, which does vary with
$\omega(b)$. The harmonic resonance gives the prime; the closure
defect gives the ring.
\end{remark}

\subsection{Continuum identification with $\sinc^2$}

\begin{theorem}[Discrete-to-continuum identity]
\label{thm:continuum}
Fix $t \in (0, 1)$ and let $f \to \infty$ along a sequence with
$k / f \to t$. Then
\begin{equation}
  R(k, f) \longrightarrow \sinc^2(t),
  \qquad \sinc(t) = \frac{\sin(\pi t)}{\pi t}.
  \label{eq:continuum}
\end{equation}
In particular, the universal mid-period constant
\[
  \lim_{f \to \infty,\, k/f \to 1/2} R(k, f) = \sinc^2(1/2) = \frac{4}{\pi^2}
  = 0.4052847\ldots
\]
\end{theorem}

\begin{proof}
By Theorem~\ref{thm:countdown},
$R(k, f) = \sin^2(\pi k / f) / (k^2 \sin^2(\pi / f))$. As
$f \to \infty$ with $k / f \to t$ fixed, we have
$\sin(\pi k / f) \to \sin(\pi t)$ and
$k \cdot \sin(\pi / f) \to k \cdot (\pi / f) = \pi t$. Hence
\[
  R(k, f)
  = \frac{\sin^2(\pi k / f)}{(k \sin(\pi / f))^2}
  \longrightarrow \frac{\sin^2(\pi t)}{(\pi t)^2}
  = \sinc^2(t).
\]
At $t = 1/2$, $\sinc^2(1/2) = (\sin(\pi/2) / (\pi/2))^2
= (2/\pi)^2 = 4/\pi^2$.
\end{proof}

\begin{remark}[Interpretation: alphabet as rectangular spectral window]
The function $\sinc^2(t)$ is the power spectral density of the
indicator $\mathbf{1}_{[0, 1]}$ on $\R$ \cite{Shannon1949,
OppenheimSchafer2010}. The alphabet $\{1, \ldots, k\}$ is a
discrete rectangular window of width $k$;
Theorem~\ref{thm:continuum} identifies the discrete spectral
power $R(k, f)$ at frequency $1/f$ with the continuous
power spectral density of the rectangular window in the
appropriate scaling limit. The discrete signal $R$ is therefore
the exact arithmetic-windowed analogue of the rectangular-pulse
power spectrum.
\end{remark}

\subsection{Stationary point at the gate}

\begin{corollary}[Stationary point at $k = p_1$]
\label{cor:stationary}
Let $f$ be a prime $\ge 3$. Then the first difference of $R$ at
$k = f$ vanishes:
\[
  R(f + 1, f) = R(f - 1, f) = \frac{1}{f^2} \cdot \frac{f^2}{(f-1)^2}
  \cdot \mathrm{(symmetry)};
\]
more precisely, $R(f-1, f) = 1/(f-1)^2$ and
$R(f+1, f) = 1/(f+1)^2$, so the symmetric difference $D_1(f) :=
R(f+1, f) - R(f-1, f) = 1/(f+1)^2 - 1/(f-1)^2 < 0$ but the
\emph{value} at $f$ is exactly $0$, i.e., $R(f, f)$ is the unique
local minimum of $R(\cdot, f)$ in the integer range
$\{f-1, f, f+1\}$.
\end{corollary}

\begin{proof}
Direct from Theorem~\ref{thm:countdown}: the values $1/(f-1)^2$,
$0$, and $1/(f+1)^2$ at $k = f-1, f, f+1$ are computed from the
closed form using $\sin^2(\pi - x) = \sin^2(x)$ and the identity
$\sin^2(\pi(f+1)/f) = \sin^2(\pi/f)$.
\end{proof}

This local-minimum structure recovers, in the integer-arithmetic
setting, the classical observation that $\sinc^2$ has a local
minimum at every positive-integer zero of $\sin$
\cite{OppenheimSchafer2010}.

%%% ==================================================================
\section{Verification}
\label{sec:verification}
%%% ==================================================================

We verified Theorems \ref{thm:countdown}--\ref{thm:omegablind}
exhaustively at machine precision. All computations use IEEE
double-precision floating point (errors at the level of
$10^{-16}$, machine epsilon).

\subsection{Closed-form identity at small primes}

For each prime $f \in \{3, 5, 7, 11, 13, 17, 19, 23\}$, we compute
$R(k, f)$ as a literal double-precision sum
$|(1/k) \sum_{j=1}^{k} e^{2\pi i j / f}|^2$ and compare to the
closed-form value $\sin^2(\pi k / f) / (k^2 \sin^2(\pi / f))$ for
all $k \in \{1, \ldots, f+1\}$. The maximum closed-form deviation
across all 168 $(f, k)$ pairs is $4.44 \times 10^{-16}$. Selected
values at the pre-collapse $k = f - 1$ exhibit the predicted
$1/(f-1)^2$:
\begin{center}
\begin{tabular}{c|c|c|c|c}
\toprule
$f$ & $k = f - 1$ & $R(f-1, f)$ measured & $1 / (f-1)^2$ predicted & error \\
\midrule
$3$  & $2$  & $0.25000000$  & $0.25000000$  & $5.55 \times 10^{-17}$ \\
$5$  & $4$  & $0.06250000$  & $0.06250000$  & $2.78 \times 10^{-17}$ \\
$7$  & $6$  & $0.02777778$  & $0.02777778$  & $3.47 \times 10^{-18}$ \\
$11$ & $10$ & $0.01000000$  & $0.01000000$  & $1.39 \times 10^{-17}$ \\
$13$ & $12$ & $0.00694444$  & $0.00694444$  & $1.91 \times 10^{-17}$ \\
$17$ & $16$ & $0.00390625$  & $0.00390625$  & $9.54 \times 10^{-18}$ \\
$19$ & $18$ & $0.00308642$  & $0.00308642$  & $1.30 \times 10^{-18}$ \\
$23$ & $22$ & $0.00206612$  & $0.00206612$  & $1.73 \times 10^{-18}$ \\
\bottomrule
\end{tabular}
\end{center}

\subsection{Macro sweep across 187 semiprimes}

For each of 187 semiprimes $b = p \times q$ with $3 \le p < q$ and
$p \le 59$ (and a representative sample of larger $q$), we
verified Theorem~\ref{thm:countdown} at $k = p - 1$ (the
pre-collapse minimum) and Theorem~\ref{thm:zerowidth}
synchronization (the partition gate event at $k = p_1$). The
macro sweep totals 36{,}662 $(b, k)$-pair evaluations of $R(k, 1/p)$;
the maximum error against the closed form is $1.11 \times 10^{-16}$
(machine epsilon). Zero exceptions in 36{,}662 cases.

Selected entries from the sweep:
\begin{center}
\begin{tabular}{c|c|c|c|c|c}
\toprule
$b$ & $p$ & $q$ & $R(p-1, 1/p)$ measured & $1/(p-1)^2$ predicted & error \\
\midrule
$35$   & $5$  & $7$  & $0.062500$ & $0.062500$ & $2.78 \times 10^{-17}$ \\
$77$   & $7$  & $11$ & $0.027778$ & $0.027778$ & $1.39 \times 10^{-17}$ \\
$143$  & $11$ & $13$ & $0.010000$ & $0.010000$ & $1.39 \times 10^{-17}$ \\
$323$  & $17$ & $19$ & $0.003906$ & $0.003906$ & $1.04 \times 10^{-17}$ \\
$667$  & $23$ & $29$ & $0.002066$ & $0.002066$ & $1.73 \times 10^{-18}$ \\
$1073$ & $29$ & $37$ & $0.001276$ & $0.001276$ & $2.39 \times 10^{-18}$ \\
$4187$ & $53$ & $79$ & $0.000370$ & $0.000370$ & $3.96 \times 10^{-18}$ \\
\bottomrule
\end{tabular}
\end{center}

\subsection{$\omega$-blindness across ring structures}

Theorem~\ref{thm:omegablind} predicts $R(k, 1/p)$ is identical for
every $b$ with $p \mid b$, regardless of $\omega(b)$. Verification
holds $p = 7$ fixed and varies $b$ across $\omega \in \{1, 2, 3\}$:
\begin{center}
\begin{tabular}{c|c|c|c}
\toprule
$b$    & $\omega(b)$ & ring structure         & $R(1, 1/7), R(2, 1/7), \ldots, R(6, 1/7)$ \\
\midrule
$49$   & $1$ & $7^2$        & $1.0000, 0.8117, 0.5610, 0.3156, 0.1299, 0.0278$ \\
$343$  & $1$ & $7^3$        & $1.0000, 0.8117, 0.5610, 0.3156, 0.1299, 0.0278$ \\
$77$   & $2$ & $7 \cdot 11$ & $1.0000, 0.8117, 0.5610, 0.3156, 0.1299, 0.0278$ \\
$91$   & $2$ & $7 \cdot 13$ & $1.0000, 0.8117, 0.5610, 0.3156, 0.1299, 0.0278$ \\
$1001$ & $3$ & $7 \cdot 11 \cdot 13$ & $1.0000, 0.8117, 0.5610, 0.3156, 0.1299, 0.0278$ \\
$1463$ & $3$ & $7 \cdot 11 \cdot 19$ & $1.0000, 0.8117, 0.5610, 0.3156, 0.1299, 0.0278$ \\
\bottomrule
\end{tabular}
\end{center}
The signal is byte-identical across all six rings (up to floating
point), as required by Theorem~\ref{thm:omegablind}.

\subsection{Universal mid-period constant}

Theorem~\ref{thm:continuum} predicts $R(k, p) \to 4/\pi^2$ as
$p \to \infty$ along $k / p \to 1/2$. Numerical verification at
proxy primes:
\begin{center}
\begin{tabular}{c|c|c}
\toprule
$p$   & $R(\lfloor p / 2 \rfloor, p)$ & deviation from $4/\pi^2 = 0.4052847\ldots$ \\
\midrule
$1009$    & $0.406090$ & $+8.0 \times 10^{-4}$ \\
$10007$   & $0.405366$ & $+8.1 \times 10^{-5}$ \\
$100003$  & $0.405293$ & $+8.1 \times 10^{-6}$ \\
\bottomrule
\end{tabular}
\end{center}
Convergence to $4/\pi^2$ is at the rate $\bigO(1/p^2)$, consistent
with the Taylor expansion of $\sin^2(\pi/p)$ around $p = \infty$.

\subsection{Software, runtime, and reproducibility}

The verification script (\texttt{verify\_first\_g.py} from the
J04 bundle, in extended-mode for J07) runs with \texttt{numpy +
sympy + math.gcd} on Python 3.10+. Total runtime under three
minutes on a 2024 consumer laptop. All sums are computed in
double-precision floating point; gcd computations use exact integer
arithmetic. The script exits 0 if and only if every comparison
falls below $10^{-10}$ (well above machine epsilon). The bundle
DOI is 10.5281/zenodo.18852047.

%%% ==================================================================
\section{Relation to Montgomery's pair correlation}
\label{sec:montgomery}
%%% ==================================================================

Montgomery \cite{Montgomery1973} proved (assuming the Riemann
hypothesis, for the relevant range of test parameters) that the
pair correlation of normalized non-trivial zeros of $\zeta(s)$
satisfies
\[
  R_2(u) = 1 - \left(\frac{\sin(\pi u)}{\pi u}\right)^2 = 1 - \sinc^2(u).
\]
Theorem~\ref{thm:continuum} of the present paper gives, in the
arithmetic side,
\[
  R(x) := \lim_{f \to \infty,\, k/f \to x} R(k, f) = \sinc^2(x).
\]
The two functions are formal complements:
\[
  R(x) + R_2(u) = \sinc^2(x) + (1 - \sinc^2(u))
  = 1 + (\sinc^2(x) - \sinc^2(u)),
\]
which equals $1$ exactly when $u = x$. At the half-period $x = u = 1/2$,
both frameworks pin the same constant:
\[
  \sinc^2(1/2) = \frac{4}{\pi^2} = 0.4052847\ldots
\]

We state the relationship as a structural identity, not as a
mechanism. The closed form for $R(k, f)$ is elementary
(geometric sum); Montgomery's $R_2(u)$ comes from a deep
analytic-number-theory argument requiring the explicit formula
and a hypothesis on $\zeta$. The complementarity is at the level
of basis functions, not at the level of consequences for the
analytic continuation of $\zeta$. We do \emph{not} claim that
Theorem~\ref{thm:continuum} bears on the Riemann hypothesis.

The numerical experiments of Odlyzko \cite{Odlyzko1987,
Odlyzko1992} confirmed Montgomery's GUE statistics to extreme
precision. The present paper's exhaustive verification across
36{,}662 semiprime test cases is the analogous empirical
confirmation on the arithmetic side. Both bodies of evidence are
consistent with the underlying $\sinc^2$ structure.

\begin{remark}[Open problem statement]
\label{rem:open-bridge}
The discrete sum $\sum_{k = 1}^{f - 1} R(k, f) \cdot g(k / f)$
for a smooth test function $g$ converges in the limit $f \to \infty$
to $\int_0^1 \sinc^2(x) g(x)\, dx$ by the Riemann-sum interpretation
of Theorem~\ref{thm:continuum}. Whether the appropriate dual sum
$\sum_{|\gamma| \le T} R_2(\gamma / 2\pi) \cdot \tilde g(\gamma)$
over Riemann zeros admits an analogous closed-form approximation
in some scaling regime is an open question we leave to future work.
\end{remark}

%%% ==================================================================
\section{Scope and limitations}
\label{sec:scope}
%%% ==================================================================

We list explicitly what the paper does and does not claim.

\begin{enumerate}[label={\upshape(\arabic*)}]
\item \textbf{What is proved.} The closed form
      $R(k, f) = \sin^2(\pi k / f) / (k^2 \sin^2(\pi / f))$
      (Theorem~\ref{thm:countdown}); the synchronization with
      the First-G partition gate at $k = p_1 = \spf(b)$
      (Theorem~\ref{thm:zerowidth}); the $\omega$-blindness
      of the harmonic signal once a prime divides $b$
      (Theorem~\ref{thm:omegablind}); the discrete-to-continuum
      identity $R \to \sinc^2$ in the appropriate scaling limit
      (Theorem~\ref{thm:continuum}).

\item \textbf{What is verified.} 8 small primes $f$ and 187
      semiprimes (36{,}662 exact computations); max deviation
      from the closed form $1.11 \times 10^{-16}$ (machine
      epsilon); zero exceptions; runtime under three minutes
      on a 2024 consumer laptop.

\item \textbf{What is \emph{not} claimed.} No statement about the
      analytic continuation of $\zeta(s)$, no consequence for
      the Riemann hypothesis (the Montgomery correspondence in
      Section~\ref{sec:montgomery} is a structural identity at
      the level of basis functions, not a mechanism). No
      polynomial-time factoring algorithm: the closed form is
      pointwise efficient, but extracting $p_1$ from $R(k, p_1)$
      still requires the alphabet to traverse $\bigO(p_1)$ steps
      before encountering the gate, which is exponential in the
      bit-length of an RSA modulus and consistent with
      sub-exponential best-known classical complexity (NFS,
      \cite{LenstraEtAl1993}). No claim about non-squarefree $b$
      beyond what the algebra already covers; the
      $\omega$-blindness statement (Theorem~\ref{thm:omegablind})
      is verified for $\omega(b) \in \{1, 2, 3\}$ and follows
      formally from the closed form for any $\omega(b) \ge 1$.

\item \textbf{Restriction to squarefree $b$ in the J-series
      verification harness.} The verification script restricts
      to squarefree $b$ because that is the case shared with the
      J04 First-G Localization companion \cite{J04Sanders} and
      with the J01 $\sigma$-rate companion \cite{J01Sanders}.
      The closed form (Theorem~\ref{thm:countdown}) and the
      $\omega$-blindness statement (Theorem~\ref{thm:omegablind})
      are not restricted to squarefree $b$; their proofs use
      only $f \in \Z_{\ge 2}$ and the geometric sum.

\item \textbf{Finite-range verification.} The exhaustive check
      is for $p \le 59$ and a representative sample of larger
      $q$. The theorems are proved in full generality; the finite
      verification is reassurance, not a pillar of the result.
\end{enumerate}

%%% ==================================================================
\section*{Acknowledgements}
%%% ==================================================================

This paper is the second in a coordinated J-series of submissions
maintained by the TIG / CK research program of 7Site LLC and the
Coherence Keeper community at
\href{https://github.com/TiredofSleep/ck}{github.com/TiredofSleep/ck}.
The First-G Localization Theorem of the companion submission
\cite{J04Sanders} is the foundational lemma on which this paper
builds. We thank Coherence Keeper collaborators for the
discrete-Fej\'er calculation that yielded
Theorem~\ref{thm:countdown} in elementary form, and
C.~A.~Luther for the dispersion conjecture that motivated the
ratio analysis underlying the seeded residue persistence framing
discussed in the corpus paper \cite{Sanders2026WP35}.

%%% ==================================================================
\begin{thebibliography}{99}
%%% ==================================================================

\bibitem{J04Sanders}
B.~R. Sanders, M.~Gish (with C.~A. Luther on the manuscript title
block), \emph{The First-G Event in the Coprimality Partition:
Stability Windows, CRT Idempotent Count, and Prime-Indexed Phase
Transitions}, submitted to \emph{Integers --- Electronic Journal of
Combinatorial Number Theory} (2026); J-series J04, 7Site LLC
research register, DOI \texttt{10.5281/zenodo.18852047}.

\bibitem{J01Sanders}
B.~R. Sanders, M.~Gish, \emph{Non-Associativity Decay in Binary
Composition Tables over $\Z/N\Z$}, submitted to \emph{Journal of
Combinatorial Theory, Series A} (2026); J-series J01, 7Site LLC
research register, DOI \texttt{10.5281/zenodo.18852047}.

\bibitem{J02Sanders}
B.~R. Sanders, M.~Gish, \emph{Joint Closure, Per-Coordinate Fuse
Data, and a Closed-Form Algebraic Attractor of Two Commutative
Binary Operations on $\Z/10\Z$}, submitted to \emph{Algebraic
Combinatorics} (2026); J-series J02, 7Site LLC research register,
DOI \texttt{10.5281/zenodo.18852047}.

\bibitem{Sanders2026WP35}
B.~R. Sanders, C.~A. Luther, M.~Gish, \emph{The Prime Phase
Transition: Harmonic Pre-Echo, Zero-Width Gates, and the Geometry
of RSA Security}, working paper WP35 (2026), 7Site LLC research
register, DOI \texttt{10.5281/zenodo.18852047}.

\bibitem{Apostol1976}
T.~M. Apostol, \emph{Introduction to Analytic Number Theory},
Undergraduate Texts in Mathematics, Springer-Verlag, New York-Heidelberg,
1976.
\textsc{MR}~\href{https://mathscinet.ams.org/mathscinet-getitem?mr=0434929}{0434929}.

\bibitem{Bochner1933}
S.~Bochner, \emph{Monotone Funktionen, Stieltjessche Integrale und
harmonische Analyse}, Math.\ Ann.\ \textbf{108} (1933), 378--410.
\textsc{MR}~\href{https://mathscinet.ams.org/mathscinet-getitem?mr=1512893}{1512893}.

\bibitem{Davenport2000}
H.~Davenport, \emph{Multiplicative Number Theory}, 3rd ed., revised
by H.~L.~Montgomery, Graduate Texts in Mathematics 74,
Springer-Verlag, New York, 2000.
\textsc{MR}~\href{https://mathscinet.ams.org/mathscinet-getitem?mr=1790423}{1790423}.

\bibitem{Fejer1900}
L.~Fej\'er, \emph{Sur les fonctions born\'ees et int\'egrables},
C.~R.~Acad.~Sci.\ Paris \textbf{131} (1900), 984--987.

\bibitem{HardyWright2008}
G.~H. Hardy and E.~M. Wright, \emph{An Introduction to the Theory
of Numbers}, 6th ed., revised by D.~R.~Heath-Brown and J.~H.~Silverman,
foreword by A.~Wiles, Oxford University Press, Oxford, 2008.
\textsc{MR}~\href{https://mathscinet.ams.org/mathscinet-getitem?mr=2445243}{2445243}.

\bibitem{IwaniecKowalski2004}
H.~Iwaniec and E.~Kowalski, \emph{Analytic Number Theory}, AMS
Colloquium Publications, Vol.\ 53, American Mathematical Society,
Providence, RI, 2004.
\textsc{MR}~\href{https://mathscinet.ams.org/mathscinet-getitem?mr=2061214}{2061214}.

\bibitem{LenstraEtAl1993}
A.~K. Lenstra, H.~W. Lenstra Jr., M.~S. Manasse, and J.~M. Pollard,
\emph{The number field sieve}, in \emph{The Development of the
Number Field Sieve}, Lecture Notes in Mathematics 1554,
Springer, Berlin, 1993, pp.~11--42.
\textsc{MR}~\href{https://mathscinet.ams.org/mathscinet-getitem?mr=1321221}{1321221}.

\bibitem{Maynard2015}
J.~Maynard, \emph{Small gaps between primes}, Ann.\ of Math.\ (2)
\textbf{181} (2015), no.~1, 383--413.
\textsc{MR}~\href{https://mathscinet.ams.org/mathscinet-getitem?mr=3272929}{3272929}.

\bibitem{Montgomery1973}
H.~L. Montgomery, \emph{The pair correlation of zeros of the zeta
function}, in \emph{Analytic Number Theory} (St.~Louis, MO, 1972),
Proc.~Sympos.~Pure Math., Vol.~XXIV, American Mathematical Society,
Providence, R.I., 1973, pp.~181--193.
\textsc{MR}~\href{https://mathscinet.ams.org/mathscinet-getitem?mr=0337821}{0337821}.

\bibitem{Odlyzko1987}
A.~M. Odlyzko, \emph{On the distribution of spacings between zeros
of the zeta function}, Math.\ Comp.\ \textbf{48} (1987), no.~177,
273--308.
\textsc{MR}~\href{https://mathscinet.ams.org/mathscinet-getitem?mr=0866115}{0866115}.

\bibitem{Odlyzko1992}
A.~M. Odlyzko, \emph{The $10^{20}$-th zero of the Riemann zeta
function and 175 million of its neighbors}, AT\&T Bell Labs preprint,
1992.

\bibitem{OppenheimSchafer2010}
A.~V. Oppenheim and R.~W. Schafer, \emph{Discrete-Time Signal
Processing}, 3rd ed., Prentice Hall, Upper Saddle River, NJ, 2010.

\bibitem{Riemann1859}
B.~Riemann, \emph{\"Uber die Anzahl der Primzahlen unter einer
gegebenen Gr\"osse}, Monatsber.\ K\"onigl.\ Preuss.\ Akad.\ Wiss.\
Berlin (Nov.~1859), 671--680.

\bibitem{Shannon1949}
C.~E. Shannon, \emph{Communication in the presence of noise},
Proc.~IRE \textbf{37}(1) (1949), 10--21.
DOI~\texttt{10.1109/JRPROC.1949.232969}.

\bibitem{Zhang2014}
Y.~Zhang, \emph{Bounded gaps between primes}, Ann.\ of Math.\ (2)
\textbf{179} (2014), no.~3, 1121--1174.
\textsc{MR}~\href{https://mathscinet.ams.org/mathscinet-getitem?mr=3171761}{3171761}.

\end{thebibliography}

\end{document}
